\chapter{计算架构与调度}

\begin{solutionexercise}{1}
\problemheading 考虑下面的 CUDA 内核及调用它的主机函数：
\begin{lstlisting}[language=C++,firstnumber=1]
__global__ void foo_kernel(int* a, int* b) {
    unsigned int i = blockIdx.x*blockDim.x + threadIdx.x;
    if (threadIdx.x < 40 || threadIdx.x >= 104) {
        b[i] = a[i] + 1;
    }
    if (i % 2 == 0) {
        a[i] = b[i] * 2;
    }
    for (unsigned int j = 0; j < 5 - (i % 3); ++j) {
        b[i] += j;
    }
}
void foo(int* a_d, int* b_d) {
    unsigned int N = 1024;
    foo_kernel<<<(N + 128 - 1)/128, 128>>>(a_d, b_d);
}
\end{lstlisting}
\begin{enumerate}[label=\alph*.,leftmargin=2.2em]
  \item 每个线程块中有多少个 warp？
  \item 网格中有多少个 warp？
  \item 对第 04 行语句：
    \begin{enumerate}[label=\roman*.,leftmargin=2em]
      \item 网格中有多少个活动 warp？
      \item 网格中有多少个分歧 warp？
      \item 块 0 的 warp 0 的 SIMD 效率（百分比）是多少？
      \item 块 0 的 warp 1 的 SIMD 效率（百分比）是多少？
      \item 块 0 的 warp 3 的 SIMD 效率（百分比）是多少？
    \end{enumerate}
  \item 对第 07 行语句：
    \begin{enumerate}[label=\roman*.,leftmargin=2em]
      \item 网格中有多少个活动 warp？
      \item 网格中有多少个分歧 warp？
      \item 块 0 的 warp 0 的 SIMD 效率（百分比）是多少？
    \end{enumerate}
  \item 对第 09 行循环：
    \begin{enumerate}[label=\roman*.,leftmargin=2em]
      \item 有多少次迭代没有分歧？
      \item 有多少次迭代产生分歧？
    \end{enumerate}
\end{enumerate}

\approachheading warp 大小为 32。先把每个 128 线程块分成四个 warp，再对每条语句
逐 warp 写出参与线程集合；SIMD 效率按“执行该语句的活动 lane 数/32”计算。

\answerheading 网格含 $\lceil1024/128\rceil=8$ 个块，每块 $128/32=4$ 个 warp，
全网格共有 $8\times4=32$ 个 warp。

对第 04 行，块内条件为 $t<40$ 或 $t\ge104$。warp 0 的 32 个线程全参与；warp 1
只有 $t=32,\ldots,39$ 的 8 个线程参与；warp 2 无线程参与；warp 3 只有
$t=104,\ldots,127$ 的 24 个线程参与。因此每块有 3 个活动 warp、2 个分歧 warp，
全网格分别为 24 和 16 个。块 0 的 warp 0、1、3 的 SIMD 效率依次为
$100\%$、$8/32=25\%$、$24/32=75\%$。

对第 07 行，连续整数索引中偶数线程执行、奇数线程不执行。32 个 warp 全都活动且
全都分歧；块 0 的 warp 0 有 16 个活动 lane，效率为 $50\%$。

第 09 行循环的迭代次数由 $i\bmod3$ 决定，分别为 5、4、3。$j=0,1,2$ 时所有线程
仍在循环中，共 3 次无分歧迭代；$j=3$ 和 $j=4$ 时只有部分线程继续，共 2 次分歧
迭代。

\conclusionheading (a) 4；(b) 32；(c) 24、16、$100\%$、$25\%$、$75\%$；
(d) 32、32、$50\%$；(e) 3 次无分歧、2 次分歧。
\discussionheading
\discussionpoint{分歧发生在动态 warp，而不是整个线程块。}
若块内 warp 0 全走真分支、warp 1 全走假分支，两组线程虽然选择不同，硬件仍可分别执行，没有 warp 内分歧；只有同一 warp 的 lane 在同一动态分支上分成两组，才需要用不同活动掩码依次处理路径。本题的区间边界恰好切过 warp 1 和 warp 3，所以它们分歧，完整落在区间内外的 warp 则一致。循环还会在每次回边重新判断参与者，静态源码只有一个 \code{for}，动态执行却产生多个不同掩码。这个层次区别也适用于变长序列和稀疏行：数据类别必须按 warp 分布分析，不能只给出全块比例。
\discussionpoint{短分支可能被谓词化，但闲置 lane 仍然存在。}
编译器有时不用显式跳转，而给几条短指令附上谓词，让 warp 发射两侧指令，仅允许满足条件的 lane 提交结果。源码级分支看似消失，活动掩码的成本却没有凭空消失：不满足谓词的 lane 仍占据发射槽，只是不改变架构状态。较长分支、嵌套循环和独立线程调度还会引入不同的重汇合行为，不能一律用“两条路径时间相加”预测。理解谓词化的意义，是知道源码分支数与机器分支数可能不同，因此最终诊断要结合生成指令和活动线程，而不是只数 \code{if}。比较编译选项前还应检查 SASS，确认谓词化边界是否真的发生变化。
\discussionpoint{重排数据能减少分歧，却可能付出更大代价。}
若条件由数据类别决定，可以先扫描并按类别分桶，让每个 warp 尽量处理同一类对象；这在一条路径包含大量计算、类别在多个阶段复用时可能值得。重排本身要读取、写入索引并可能进行全局同步，若分支体只有一次加法，它的成本往往比原分歧更高。类别数量多或频繁变化时，分桶还会破坏原有空间局部性和稳定顺序，需要额外元数据恢复结果。因而优化问题不是“能否消灭分歧”，而是“节省的被屏蔽指令是否超过分类、搬运与顺序恢复的总成本”。输入分布改变时还要重新评估，因为稳定训练集上的桶比例未必代表线上请求。
\discussionpoint{诊断需要构造具有不同活动掩码的输入。}
只用均匀随机数据测一次，很难判断慢是由分支、访存还是尾部造成。更有解释力的实验让条件分别全真、全假、lane 交替和按半 warp 分组，再比较每条源码的动态实例、平均活动 lane 和总时间；这些输入具有相同规模，却形成截然不同的掩码。若全真与全假都快而交替明显变慢，分歧解释才获得支持；若四者接近，瓶颈可能在内存或另一阶段。测量应围绕可证伪的假设设计，而不是堆叠一串指标后凭相关性下结论。还应固定每条路径的数据地址，避免掩码实验同时改变缓存命中而混淆因果。
\variantheading 若第 3 行条件改为 $t<32$ 或 $t\ge96$，每块有 2 个活动 warp且都不分歧，全网格为 16 个活动 warp、0 个分歧 warp。
\practiceheading 用 Nsight Compute 的 Branch Efficiency、Warp State Statistics 和源码关联视图验证动态分支，而不要仅从源码推断实际指令。
\sourcecheck{https://docs.nvidia.com/cuda/cuda-c-programming-guide/}{NVIDIA CUDA C++ Programming Guide 对 32 线程 warp、活动线程与分歧的定义；逐 lane 枚举复核，置信度高}
\end{solutionexercise}

\begin{solutionexercise}{2}
\problemheading 对向量加法，假设向量长度为 2000，每个线程计算一个输出元素，线程块
大小为 512。网格中有多少个线程？

\approachheading 块数必须向上取整，实际启动线程数等于块数乘块大小。

\answerheading
\[
 B=\left\lceil\frac{2000}{512}\right\rceil=4,
 \qquad T_{\rm launch}=4\times512=2048\ \text{个线程}.
\]
其中 2000 个线程处理元素，48 个线程由边界检查屏蔽。

\conclusionheading 网格启动 2048 个线程（4 个线程块）。
\discussionheading
\discussionpoint{任意数组长度必须适配固定的执行粒度。}
数组可以有 2000 个元素，硬件却按完整线程块分派、按完整 warp 发射，因此四个 512 线程块覆盖到索引 2047 是正常的离散化结果。边界谓词只保证最后 48 个线程不产生数组副作用，并不会让硬件把它们从块中删掉；它们仍执行索引和条件判断。这个现象与存储系统按整页分配、网络按整包传输相似：逻辑有效载荷不必填满物理容器。评估冗余时应区分必须支付的容器开销和可以通过改变任务组织减少的重复工作。这里的 48 个空线程只占一个块尾部，不能据此推断整个网格都以同样比例浪费。
\discussionpoint{专用尾内核只有在尾部足够重要时才划算。}
理论上可以让主体内核只处理完整块，再启动一个更小的内核计算最后 464 项，从而减少 48 个空线程；但第二次启动、参数设置和代码维护通常远贵于几十次条件判断。若程序连续处理数百万个短向量，或每个尾元素执行昂贵且分歧严重的算法，把尾部批量收集到专门内核才可能收回成本。另一选择是融合多个小数组，让一个块处理一批对象，同时减少启动次数与空槽。优化应比较端到端时间，而不是以“没有空线程”为目标，因为消除局部浪费可能引入更大的全局开销。专用路径还会扩大测试矩阵，必须覆盖主体与尾部交界处的不重不漏。
\discussionpoint{块内尾线程与网格最后一波是两个独立问题。}
48 个无效线程位于最后一个块内部，它们影响该块的活动 lane；“最后一波不足”则取决于整个网格有多少块相对所有 SM 的并发块容量。即使 $N$ 恰好整除 512、块内没有一个尾线程，只有 4 个块的网格在拥有许多 SM 的设备上仍可能从一开始就填不满；反过来，长网格有数百波时，最后一个不满块的固定成本很容易摊薄。因此不能用 $48/2048$ 推导设备利用率，也不能把小网格尾延迟归因于这 48 个线程。前者可通过块宽改善，后者需要更多独立块或任务融合，修复手段并不相同。
\discussionpoint{块大小应在真实长度分布上选择。}
把块从 512 改成 256，本题仍启动 2048 个线程，却会形成 8 个块，给调度器更细粒度的分配机会，同时改变每块寄存器与共享内存总量。更小并不总更好：块级协作空间会缩小，固定开销增加，某些归约还需要更多阶段。可以先用资源模型排除无法获得合理驻留数的配置，再在业务实际出现的短、中、长向量上测量中位和高分位延迟。这样选出的块宽反映长度分布与目标架构，而不是只对 $N=2000$ 这个单点偶然合适。若网格本来很小，还应单独记录 SM 覆盖率，防止较大块造成并行块数量不足。
\variantheading 若长度为 2049、块大小仍为 512，则需 5 个块、2560 个线程，其中 511 个线程越界。
\practiceheading 在生产配置中用 \code{cudaOccupancyMaxPotentialBlockSize} 给出候选块大小，再通过 Nsight Compute 比较占用率与尾块浪费。
\sourcecheck{https://docs.nvidia.com/cuda/cuda-c-programming-guide/}{NVIDIA CUDA C++ Programming Guide 的线程层次与网格维度定义；算式独立复核，置信度高}
\end{solutionexercise}

\begin{solutionexercise}{3}
\problemheading 对上一题，考虑长度为 2000 的边界检查，预计有多少个 warp 会产生分歧？

\approachheading 按 32 个连续索引一组检查最后的有效/无效线程分界。

\answerheading $2000=62\times32+16$。前 62 个 warp 全部有效；下一个 warp 中
16 个线程有效、16 个线程越界，因而发生分歧；最后一个 warp 的 32 个线程全部越界，
整个 warp 走同一分支，不发生分歧。

\conclusionheading 只有 1 个 warp 因边界检查发生分歧。
\discussionheading
\discussionpoint{没有分歧的 warp 也可能完全没有有效工作。}
$2000=62\times32+16$，第 63 个 warp 有 16 个有效 lane 和 16 个越界 lane，因而在边界分支处分歧；再后的 warp 若全部越界，会一致跳过主体，分支效率反而可能是百分之百。后者仍被创建、调度并执行条件，只是没有 lane 写输出，所以“零分歧”不能推出“利用率高”。同理，一个全真分支可以无分歧却执行昂贵路径，一个半真分支也可能因主体极短而成本很小。性能语言必须把控制一致性与有用工作比例分开，否则优化会追逐错误目标。分析器中的活动线程指标与输出计数应同时使用，才能识别这种全空 warp。
\discussionpoint{余数决定部分 warp，块取整决定全空 warp。}
若按连续索引映射，$N\bmod32$ 给出最后一个部分 warp 的有效 lane 数；当余数为零时，边界恰落在 warp 之间，部分 warp 数为零。块大小再向上取整到完整线程块，可能在其后留下若干完全空 warp，本题 512 线程块的最后一块便同时含有部分和空闲区域。分别统计有效线程、部分 warp 和全空 warp，能精确解释边界判断的动态行为。这个计数还可预测改变块宽后哪些成本会消失，而一项总体“分支效率”无法提供同样信息。若采用二维映射，余数应沿每一维分别计算，角块不能用一维公式直接代替。
\discussionpoint{为了消除一个尾分支而增加第二条路径常得不偿失。}
可以让主体只覆盖 $\lfloor N/32\rfloor\times32$ 个元素，再由一个 warp、另一个内核甚至 CPU 处理余数，从而使主体完全一致。代价是多一套调度与边界代码，若尾部只有一次加法，额外启动远大于 16 个被屏蔽 lane 的损失。带掩码的库原语可保持一条调用路径，但底层仍必须确保无效 lane 不读写越界，并没有消除物理空闲。只有尾部分支很长、被大量短数组反复触发，或能把许多尾部合并处理时，专门化才具有系统层面的收益。决定拆分前可建立阈值实验，找出尾部成本何时超过额外调度与同步成本。
\discussionpoint{变长数据把一个尾部扩散到每个样本。}
ragged batch、稀疏矩阵行和自然语言序列都具有不同长度，尾部不再只出现于整个数组末端，而可能在每个 warp 负责的样本中反复出现。按长度分桶可让相近样本共用 warp，padding 则用额外存储换规则控制流，压缩索引还能只调度有效项；三者分别支付排序、填充流量和间接寻址成本。最佳选择取决于长度分布是否稳定、后续算子能否复用分桶以及填充值是否参与昂贵计算。这个更广场景说明，本题的一次余数计算其实是负载整形问题的最小模型。线上长度分布若漂移，固定桶边界也应重新测量，而不是永久沿用离线最优值。
\variantheading 若向量长度为 1984，即 $62\times32$，边界落在 warp 边界上，分歧 warp 为 0，最后两个 warp 全部越界。
\practiceheading Nsight Compute 的 Branch Efficiency 应结合 Active Threads Per Warp 和 Executed Ipc 分析，避免把“无分歧但全被屏蔽”误判为高效。
\sourcecheck{https://docs.nvidia.com/cuda/cuda-c-programming-guide/}{NVIDIA CUDA C++ Programming Guide 的 SIMT 分支行为；边界余数独立复核，置信度高}
\end{solutionexercise}

\begin{solutionexercise}{4}
\problemheading 假设一个有 8 个线程的线程块在到达屏障前执行一段代码。8 个线程执行
这段代码所需时间（微秒）分别为 2.0、2.3、3.0、2.8、2.4、1.9、2.6 和
2.9，其余时间都在等待屏障。线程总执行时间中有百分之多少用于等待屏障？

\approachheading 屏障释放时间由最慢线程的 $3.0\ \mu\mathrm{s}$ 决定；把每个线程
相对该时刻的空闲时间相加，再除以 8 个线程的总驻留时间。

\answerheading 有效执行时间之和为
\[
2.0+2.3+3.0+2.8+2.4+1.9+2.6+2.9=19.9\ \mu\mathrm{s}.
\]
到屏障释放为止的总线程时间为 $8\times3.0=24.0\ \mu\mathrm{s}$，故等待时间为
$24.0-19.9=4.1\ \mu\mathrm{s}$，比例为
\[
 \frac{4.1}{24.0}\times100\%=17.0833\%\approx17.1\%.
\]

\conclusionheading 约 $17.1\%$ 的线程总执行时间用于等待屏障。
\discussionheading
\discussionpoint{把等待量看成关键路径周围的“松弛”。}
为什么八个线程都做了有用工作，仍能累计出 $4.1\ \mu\mathrm{s}$ 的等待？在题设的同时开始模型中，屏障释放时刻是 $t_{\max}=3.0\ \mu\mathrm{s}$，线程 $i$ 的松弛量为 $t_{\max}-t_i$，所以等待比例一般写成 $\sum_i(t_{\max}-t_i)/(n t_{\max})$。这个比值既取决于最慢者，也取决于其余时间分布；仅把最快线程再加速，甚至可能让总等待比例上升而不缩短关键路径。它与并行任务图中的关键路径和工作量下界相呼应，但结论依赖“同时开始、持续占用且时间可相加”的抽象，不能直接当作真实 GPU 的能耗或利用率。

\discussionpoint{单线程计时怎样映射到 warp 调度。}
真实 GPU 不会给八个 lane 各配一个完全独立的指令调度器，因此题中的八个时间不能机械解释成八条硬件时间线。同一 warp 内的差异通常来自控制流分歧、不同长度的依赖链或不同访存完成时刻，而调度器发射的是 warp 指令；若这些线程分属多个 warp，其他就绪 warp 又可能在等待期间发射，从而隐藏一部分延迟。屏障自身还有到达速率和同步指令吞吐量，最后一个 warp 何时到达才决定释放。因此 $17.1\%$ 是给定教学模型下的线程时间占比，不是分析器中 SM 利用率、占用率或“stall barrier”百分比的同义词。

\discussionpoint{减少等待不能等同于删除屏障。}
若慢线程处理高次数循环而快线程只处理短任务，首先应问工作能否重新分桶、用动态队列分配，或让每个 warp 承担更均匀的批次；这些做法缩小到达时间分布，才真正缩短屏障前关键路径。某些算法还能让装载 warp 与计算 warp 专门化，或把相邻批次流水化，使等待与独立工作重叠。可是只要下一阶段会读取所有生产者写入的数据，直接删掉屏障就破坏到达条件与内存可见性，性能问题会转化成数据竞争。更一般的设计原则是先画出生产者--消费者依赖，再缩小同步作用域、减少同步次数或均衡作用域内工作，而不是凭一次计时猜测屏障“多余”。

\discussionpoint{从症状反推真正的慢到达者。}
分析器报告某 warp 因 barrier 停顿，只说明它已经到达而同组参与者尚未满足条件，并不告诉我们迟到者是在算术、分支还是内存依赖上耗时。较可靠的实验会分别关闭或替换候选阶段，记录阶段级 GPU 时间戳，并比较线程块之间与块内 warp 之间的到达分布；若短阶段被时间戳写入本身明显扰动，还应改用多个独立内核和事件做差分。随后把 barrier stall 与分支效率、内存延迟、eligible warps 和调度器发射情况一起看，才能判断均衡工作是否有效。这样的诊断也适用于 CPU 线程池和分布式栅栏：等待者容易观测，真正需要优化的却通常是尚未到达的路径。
\variantheading 若最慢线程由 $3.0$ 降至 $2.9\ \mu\mathrm{s}$，有效执行时间之和变为 $19.8\ \mu\mathrm{s}$，总线程时间为 $23.2\ \mu\mathrm{s}$，等待为 $3.4\ \mu\mathrm{s}$，比例约 $14.7\%$。
\practiceheading 用 Nsight Compute 的 Barrier、Warp Stall Barrier 和 Scheduler Statistics 指标区分屏障等待与内存依赖停顿。
\sourcecheck{https://docs.nvidia.com/cuda/cuda-c-programming-guide/}{NVIDIA CUDA C++ Programming Guide 对块级屏障的语义；时间总量独立求和，置信度高}
\end{solutionexercise}

\begin{solutionexercise}{5}
\problemheading CUDA 程序员认为，如果每个线程块只启动 32 个线程，那么在需要屏障同步
的地方可以省略 \code{__syncthreads()}。你认为这是好主意吗？请解释原因。

\approachheading 区分“线程属于同一 warp”和“程序获得了内存可见性及 happens-before
保证”。前者不能自动替代后者。

\answerheading 这不是可靠做法。自 Volta 起，独立线程调度允许同一 warp 的线程以
更细粒度暂停和推进；依赖隐式 lockstep 的 warp-synchronous 代码可能读到尚未写入的
共享数据。即使在较早架构上，编译器重排和内存可见性也不应靠未写入程序的调度假设
保证。若通信确实只涉及一个 warp，可用带正确参与掩码的 \code{__syncwarp(mask)}；
若语义要求整个线程块到齐并让先前共享/全局内存访问对块内线程可见，就必须使用
\code{__syncthreads()}。所有参与线程还必须以一致控制流到达块级屏障。

\conclusionheading 不能仅因块大小为 32 就删去同步。应根据所需作用域选用 warp
同步或块级同步。
\discussionheading
\discussionpoint{锁步直觉缺少哪一半语义。}
假设 lane 0 写共享变量而 lane 1 随后读取，即使人眼觉得一个 warp 会“整齐向前”，程序仍需要同时回答两个问题：消费者是否等到生产者到达，以及写入是否按规定对消费者可见。屏障把到达同步与相应作用域内的内存顺序组合起来，未写入程序的调度习惯并不能替代这份契约。自 Volta 起的独立线程调度使同一 warp 中不同 lane 可以更细粒度地暂停和推进，让旧式 warp-synchronous 技巧更容易暴露；但即使面向更早设备，也不应依赖编译器重排和硬件时序的偶然结果。把同步写明还能使代码审查围绕 happens-before 关系进行，而不是围绕某一代 GPU 的经验猜测。

\discussionpoint{同步原语应覆盖恰好需要通信的参与者。}
若数据只在掩码指定的一个 warp 内交换，\code{__syncwarp(mask)} 可以等待这些 lane 并提供规定的内存排序；若生产者和消费者跨越多个 warp，块级 \code{__syncthreads()} 才覆盖完整通信域。再扩大到线程块集群、整个网格或不同内核时，需要 cooperative groups、受支持的集群同步、协作启动，或以内核边界和事件建立顺序，不能把块级屏障想象成全设备栅栏。作用域越大，必须同时到达的工作越多，慢参与者带来的等待也越广。因而选择“最小但充分”的作用域既是正确性问题，也是可扩展性问题，这与 CPU 原子操作中的内存序和分布式系统中的集合通信具有相同结构。

\discussionpoint{边界线程提前退出为何会破坏整个块。}
设一个二维块覆盖图像右下角，只有部分线程对应有效像素；若无效线程在共享装载前直接 \code{return}，剩余线程随后执行块级屏障，参与集合就不再一致，可能挂起或产生未定义行为。安全结构是让所有线程用谓词装载有效输入或零，所有线程无条件到达屏障，然后只让有效线程计算和写回。类似地，把 \code{__syncthreads()} 放进仅部分线程成立的普通分支也不安全，除非能证明条件对整个块一致。这个反例说明边界检查的位置属于同步证明的一部分：谓词可以屏蔽数据操作，却不能随意改变集体操作的成员集合。

\discussionpoint{动态检查只能辅助依赖证明。}
Compute Sanitizer 的 racecheck 能暴露许多共享内存冲突，synccheck 能发现一部分非法屏障或 warp 掩码使用，但一次测试通过并不证明所有输入和调度下都存在正确的 happens-before。更稳妥的审查会为每个共享数组列出写者、读者、索引域和跨线程边，再在每条边上标注负责排序的同步原语；如果某条边跨块，块级屏障立即显得不充分。还应构造尾块、分支两侧和不同 warp 到达次序等对抗用例，而不是只运行整齐尺寸。这种“依赖图加动态检测”的组合也能迁移到 CPU 并发代码：工具寻找实际竞态，结构化证明覆盖尚未发生的调度。
\variantheading 若通信只发生在 lane 0--15，可用参与掩码 \code{0x0000ffff} 调用 \code{__syncwarp}，前提是这 16 个 lane 都会到达该调用。
\practiceheading 采用 Compute Sanitizer racecheck/synccheck 检查共享内存竞态和非法屏障参与，并为 Volta 及更新架构编译测试独立线程调度路径。
\sourcecheck{https://docs.nvidia.com/cuda/cuda-c-programming-guide/}{NVIDIA CUDA C++ Programming Guide 的独立线程调度与同步函数语义；对抗审查确认隐式 warp 同步不是可移植保证}
\end{solutionexercise}

\begin{solutionexercise}{6}
\problemheading 如果 CUDA 设备的 SM（流式多处理器）最多容纳 1536 个线程和 4 个线程块，
下列哪一种线程块配置会使 SM 中的线程数最多？
\begin{enumerate}[label=\alph*.,leftmargin=2.2em]
  \item 每块 128 个线程；
  \item 每块 256 个线程；
  \item 每块 512 个线程；
  \item 每块 1024 个线程。
\end{enumerate}

\approachheading 对块大小 $T$，驻留块数为 $\min(4,\lfloor1536/T\rfloor)$。

\answerheading
\begin{center}
\begin{tabular}{c@{\qquad}c@{\qquad}c}
\toprule
每块线程数 & 驻留块数 & 驻留线程数\\
\midrule
128  & 4 & 512\\
256  & 4 & 1024\\
512  & 3 & 1536\\
1024 & 1 & 1024\\
\bottomrule
\end{tabular}
\end{center}

\conclusionheading 配置 (c)，每块 512 个线程，使 SM 中驻留线程数最多，为 1536。
\discussionheading
\discussionpoint{占用率首先是一个整数装箱问题。}
为什么 512 线程块能装满 1536 个线程，而 1024 线程块反而只能驻留 1024 个？对题中只考虑线程数与块数的模型，块大小为 $T$ 时驻留块数是 $B=\min(4,\lfloor1536/T\rfloor)$，驻留线程数才是 $BT$；资源不能按半个块分配，所以每次向下取整都会留下空隙。由此得到的是阶梯函数而非连续曲线，块大小略变就可能跨过一个整数边界。这个现象与内存分页、箱装和批处理中的内部碎片相似，也解释了为何“块越大越接近上限”的直觉会在 1024 线程处失效。

\discussionpoint{题目省略的资源会怎样改变答案。}
手算中 512 线程块能驻留三块，是因为我们按题意假定寄存器、共享内存和 warp 数都不限制；真实内核却要同时满足每 SM 寄存器总量、每块静态与动态共享内存、驻留 warp 数以及架构规定的分配粒度。若每块共享内存用量使 SM 只能放两块，同一配置便从 1536 线程降到 1024；若寄存器按 warp 粒度向上取整，简单的“每线程乘线程数”也可能偏乐观。编译器版本、目标 SM、循环展开和模板参数都会改变这些资源。因此这个答案是题给抽象下的确定结论，而不是同一源代码跨设备都能获得的常数。

\discussionpoint{更多驻留线程为何不保证更快。}
高占用率的主要价值是提供更多就绪 warp，使一个 warp 等待内存或依赖时调度器仍有工作可发射；它不是计算单元利用率或程序吞吐量本身。一个算术密集、指令级并行充分的内核可能用较少 warp 就填满管线，而满占用的内核若访问未合并、分支严重或每条指令都有长依赖，仍会很慢。块过小还可能削弱共享瓦片复用、归约树效率或特定矩阵指令的组织方式。因而“512 线程使驻留线程最多”只回答资源容量问题，性能选择还要同时看有效带宽、发射槽利用和每个输出的工作量。

\discussionpoint{把手算结果变成可解释的调优起点。}
工程上可以先用 Occupancy API 结合真实函数属性，排除因寄存器或动态共享内存而不可行的块大小，再在 128、256、512 等候选间进行基准测试。测试不应只选一个整齐大数组，还要覆盖生产中的尺寸分布，因为末波块数、缓存工作集和短内核启动成本都会改变排名；记录时间时同时保存寄存器、共享内存、实际驻留块和网格波次数，才能解释结果。成熟库常把最佳配置按架构、数据类型和形状缓存，而不是宣称一个全局最佳数字。这样保留了题目整数模型的预测力，也允许测量指出另一个资源或算法瓶颈已经成为主导。
\variantheading 若块上限提高到 8 而线程上限仍为 1536，则 256 线程/块可驻留 6 块并达到 1536 线程，与 512 线程/块并列最优。
\practiceheading 用 CUDA Occupancy API 或 Nsight Compute Occupancy 面板输入真实寄存器和动态共享内存用量，避免仅按线程数选择块大小。
\sourcecheck{https://docs.nvidia.com/nsight-compute/NsightCompute/index.html}{NVIDIA Nsight Compute Occupancy Calculator 对线程与块资源上限的联合约束；逐项整数除法复核，置信度高}
\end{solutionexercise}

\begin{solutionexercise}{7}
\problemheading 假设设备每个 SM 最多允许 64 个线程块和 2048 个线程。指出下列每种
每个 SM 的分配是否可行；若可行，给出占用率：
\begin{enumerate}[label=\alph*.,leftmargin=2.2em]
  \item 8 个块，每块 128 个线程；
  \item 16 个块，每块 64 个线程；
  \item 32 个块，每块 32 个线程；
  \item 64 个块，每块 32 个线程；
  \item 32 个块，每块 64 个线程。
\end{enumerate}

\approachheading 题目只给出块数和线程数限制，故检查两项上限，并以驻留线程数除以
2048。这里不额外假设寄存器、共享内存或架构分配粒度。

\answerheading
\begin{center}
\begin{tabular}{c@{\quad}c@{\quad}c@{\quad}c}
\toprule
小题 & 分配 & 驻留线程 & 可行性与占用率\\
\midrule
(a) & $8\times128$  & 1024 & 可行，$50\%$\\
(b) & $16\times64$  & 1024 & 可行，$50\%$\\
(c) & $32\times32$  & 1024 & 可行，$50\%$\\
(d) & $64\times32$  & 2048 & 可行，$100\%$\\
(e) & $32\times64$  & 2048 & 可行，$100\%$\\
\bottomrule
\end{tabular}
\end{center}

\conclusionheading 五种分配都可行；占用率依次为 $50\%,50\%,50\%,100\%,100\%$。
\discussionheading
\discussionpoint{同样的驻留线程数可以代表不同的调度结构。}
配置 (d) 与 (e) 都放入 2048 个线程并得到题设意义下的 $100\%$ 占用率，但前者是 64 个 32 线程块，后者是 32 个 64 线程块；相同总数并没有抹去块边界。较多小块给调度器更多独立完成与替换的单位，却让每块只能在较小范围内共享数据和同步；较大块能组织跨 warp 的归约或瓦片，也可能让一个慢块占据资源更久。若网格总块数很少，两种配置的末波利用率还会不同。因此占用率是一个容量投影，不是对块级协作、负载均衡和调度尾部的完整描述。比较两者时应保持算法工作相同，避免把块形状改变引起的复用差异误算成调度收益。

\discussionpoint{资源分配粒度会制造看不见的碎片。}
题目允许把 64 个单 warp 块精确装入 SM，是因为只给了线程数和块数上限；实际设备往往按 warp 或块的粒度分配寄存器、共享内存和调度状态。假如每块哪怕保留一份固定共享结构，64 个小块的固定开销就可能先触及块或内存限制，而较大块能摊薄它；反过来，单个大块的资源向上取整也可能留下无法再容纳另一块的空洞。动态共享内存和模板特化还会在运行或编译配置间改变这个结果。故比较等占用配置时必须重新读取内核属性并做多资源的整数装箱，不能只把总线程数相除。实际分配粒度应查目标架构资料或工具输出，不能从一个占用率结果反向猜定。

\discussionpoint{理论占用率与执行利用率回答不同问题。}
理论占用率表示可驻留 warp 数相对硬件上限的比例，却没有保证这些 warp 在某周期已就绪，更没有保证 ALU、加载存储或特殊函数管线都在做有用工作。长依赖链会使许多驻留 warp 同时等待，bank 冲突和缓存未命中会延长数据到达，分支分歧则让一次发射只有部分 lane 有效；于是 $100\%$ 占用率仍可能伴随很低吞吐。反之，只要较少 warp 已足以覆盖延迟，$50\%$ 也可能接近峰值。分析时应把 achieved occupancy、eligible warps、发射槽和各管线吞吐联系起来，而不是把一个百分数解释为“GPU 忙碌程度”。

\discussionpoint{块形状还编码了算法的数据几何。}
在矩阵转置中，同样 256 个线程可以排成 $32\times8$ 或 $16\times16$，两者对连续内存方向、共享存储布局和边界浪费的影响不同；在模板计算中，块边长还决定 halo 与有效输出的比例。归约需要足够线程形成树，而每线程多元素又会改变寄存器和访存数量，所以只比较线程总数会遗漏真正的算法代价。生产调优通常把二维或三维形状、每线程工作量、共享瓦片和资源上限作为联合参数，并用真实的非方形尺寸分布评价。题中五项计算教会我们先筛掉非法点，后续性能模型则要把任务几何重新放回来。
\variantheading 若设备块上限改为 32，则配置 (d) 的 64 块不可行，其余四项仍可行，(e) 仍以 2048 线程达到 $100\%$。
\practiceheading 在 Nsight Compute 中同时观察 Theoretical/ achieved Occupancy、Eligible Warps 和 SM Throughput，确认提高驻留数是否真正改善延迟隐藏。
\sourcecheck{https://docs.nvidia.com/nsight-compute/NsightCompute/index.html}{NVIDIA Nsight Compute Occupancy Calculator 的占用率定义；只按题给资源进行对抗复算，置信度高}
\end{solutionexercise}

\begin{solutionexercise}{8}
\problemheading 某 GPU 的硬件限制为：每个 SM 2048 个线程、32 个线程块和 64 K（65,536）
个寄存器。对下列内核特征，判断内核能否达到满占用率；如果不能，指出限制因素：
\begin{enumerate}[label=\alph*.,leftmargin=2.2em]
  \item 每块 128 个线程，每线程 30 个寄存器；
  \item 每块 32 个线程，每线程 29 个寄存器；
  \item 每块 256 个线程，每线程 34 个寄存器。
\end{enumerate}

\approachheading 先求容纳 2048 个线程所需块数，再核对块数和寄存器总量。按题目
抽象模型使用“每线程寄存器数乘线程数”，不引入未给出的架构分配粒度。

\answerheading
\begin{enumerate}[label=\alph*.,leftmargin=2.2em]
\item 需要 $2048/128=16$ 个块，寄存器数为
$2048\times30=61{,}440<65{,}536$，且 $16<32$，可以达到满占用率。
\item 需要 $2048/32=64$ 个块，超过 32 块上限。最多驻留
$32\times32=1024$ 个线程，即 $50\%$；限制因素是每 SM 块数。
\item 需要 $2048/256=8$ 个块，但寄存器需求
$2048\times34=69{,}632>65{,}536$。每块用 $256\times34=8704$ 个寄存器，
最多驻留 $\lfloor65536/8704\rfloor=7$ 个块，即 1792 个线程、$87.5\%$；限制因素
是寄存器。
\end{enumerate}

实际 GPU 还按架构规定的粒度分配寄存器，实测占用率可能更低，应以目标设备属性和
占用率计算器为准。

\conclusionheading (a) 可以；(b) 受块数限制；(c) 受寄存器限制。
\discussionheading
\discussionpoint{简化的寄存器方程为何有用又为何不够。}
在题设模型中，只要把每线程寄存器数 $r$ 与目标驻留线程数 $T$ 相乘，就能看出 (c) 所需 $2048\times34=69{,}632$ 已超过 65,536；这个容量反例说明线程上限允许并不代表资源组合可行。真实 GPU 却常按 warp 或块的寄存器分配单元向上取整，所以每块实际消耗可能大于 $r$ 与线程数的朴素乘积，临界的 (a) 也应交给目标架构的 occupancy 计算器复核。教学方程适合揭示限制因素和给出上界，分配粒度则解释为何实测驻留块数有时低于手算值。两层模型并不矛盾，而是分别服务于推理与部署。

\discussionpoint{寄存器数是编译产物而非源码常量。}
同一段 CUDA C++ 在循环展开、函数内联、模板实例化或目标 SM 改变后，变量的活跃区间和指令选择都会变化，因而 \code{ptxas -v} 报告的寄存器数只能归属于某个具体构建。程序员有时用 \code{--maxrregcount} 强迫编译器降低用量，希望多驻留一个块；若被挤出的值发生 spill，它们会进入 local address space，并可能增加缓存和显存流量，最终比低占用更慢。还要区分编译器报告的 local memory、stack 和 spill load/store，不能看到“local”就断言都是寄存器溢出。正确做法是保存构建参数与资源报告，再用运行时间验证这一交换是否真的获利。

\discussionpoint{算法复用与寄存器压力常来自同一个选择。}
让一个线程计算多个输出，可以把已加载的数据或系数重复使用，并暴露相互独立的乘加来提高指令级并行；代价是多个累加器和预取值同时活跃，寄存器数随粗化因子增长。软件流水线也会在当前数据计算时保留下一批数据，具有同样矛盾。若寄存器增加使驻留 warp 跨过一个装箱台阶，延迟隐藏可能突然下降；若复用减少了更多显存访问，较低占用仍可能更快。因此优化目标不是把每线程寄存器数压到最小，而是在数据复用、独立指令和可驻留并行度之间寻找可测量的平衡。扫描粗化因子时应同时记录 spill，防止把寄存器压力转移成隐藏的设备内存流量。

\discussionpoint{如何把资源跃迁纳入性能回归。}
一次编译器升级可能让模板实例多用两个寄存器，恰好使驻留块数从八个降为七个，即使源码 diff 看不出性能风险。可靠的回归应按目标架构保存寄存器数、静态与动态共享内存、stack/local 用量、spill 指令、理论和实际占用率以及内核时间，并用代表性尺寸识别是资源跃迁还是缓存变化。对靠近硬件上限的内核，还可在启动前查询函数属性并检查动态共享内存设置，避免换到资源较小设备才得到配置错误。这样的记录把题中的一次手算扩展为可追踪的构建契约，也让“更高占用却更慢”的结果有证据可解释。
\variantheading 若 (c) 每线程降为 32 个寄存器，则 8 块共需 65,536 个寄存器，可在题给简化模型下达到满占用率。
\practiceheading 用 \code{nvcc --ptxas-options=-v} 查看静态寄存器数，再调用 CUDA
Occupancy API，并用 Nsight Compute 核对实际驻留块数。
\sourcecheck{https://docs.nvidia.com/nsight-compute/NsightCompute/index.html}{NVIDIA Nsight Compute Occupancy Calculator 对寄存器、块和线程限制的说明；按题给简化模型复核，置信度高}
\end{solutionexercise}

\begin{solutionexercise}{9}
\problemheading 一名学生说，他用每块 $32\times32$ 个线程的矩阵乘法内核成功计算了两个
$1024\times1024$ 矩阵的乘法。该学生使用的 CUDA 设备每块最多允许 512 个线程、
每个 SM 最多允许 8 个线程块；并且每个线程计算结果矩阵的一个元素。你对此有
什么反应？为什么？

\approachheading 首先检查启动配置是否合法，而不是从矩阵维度推断程序已经正确执行。

\answerheading $32\times32=1024$ 个线程，超过设备每块 512 个线程的硬上限。因此
该内核启动应失败，不能完成所声称的计算；若程序没有紧接启动调用
\code{cudaGetLastError()} 并在适当位置同步检查异步错误，旧输出或未初始化输出可能
被误判为成功。每 SM 最多 8 个块并不能放宽单块线程上限。

可改用 $32\times16$ 或 $16\times16$ 的块，并把网格两维分别设为覆盖 1024 行、
1024 列的向上取整值；每个线程仍可计算一个结果元素。

\conclusionheading 该说法与硬件限制矛盾；原配置不可能合法启动，必须缩小线程块并
检查 CUDA 错误。
\discussionheading
\discussionpoint{先区分“不能启动”与“启动后可能较慢”。}
$32\times32=1024$ 超过题设每块 512 线程，是启动配置的硬约束：运行时不能把这个块拆成两个合法块，也不会以较低占用率勉强执行。每 SM 最多八块描述的是若干已经合法的块能否并发驻留，它与每块线程数、单维尺寸和每块共享内存上限处在不同层次，不能互相抵消。相反，寄存器较多而仍能驻留一块通常属于性能或资源装箱问题，不必然使启动非法。审查 CUDA 程序时先按“设备级、SM 级、块级”分层列约束，可以避免把软性能建议误当错误，也避免把真正的非法配置包装成占用率问题。

\discussionpoint{异步执行怎样制造“我算成功了”的假象。}
内核启动通常只把工作排入队列，主机若不立即检查 \code{cudaGetLastError()}，又不在读取结果前检查同步 API 的返回值，非法配置可能直到很晚才被发现。此时程序可能打印先前运行残留的缓冲区、初始化常量或根本未更新的数据；若测试矩阵全零、单位阵或输出未清空，错误结果甚至会碰巧像正确答案。可靠证据必须同时包含启动状态、完成状态和与独立 CPU 参考的逐元素比较，并让负测试确认非法配置确实被测试框架捕获。这个原则也适用于其他异步系统：提交成功不等于任务成功，更不等于结果正确。

\discussionpoint{合法块形状只是矩阵乘法重构的第一步。}
把块改为 $32\times16$ 会恰好满足 512 线程上限，改为 $16\times16$ 则给资源和调度更多余量；两者都必须重新计算网格覆盖和行列边界，才能保持每线程一个输出的语义。若追求性能，成熟 GEMM 往往让线程协作加载共享瓦片并让每线程累加多个输出，这会改变合并访问方向、数据复用、寄存器累加器和同步次数，而非简单缩小二维尺寸。还要检查矩阵维度不是块尺寸整数倍的尾块，以及 leading dimension、转置和批处理布局。由非法配置出发，可以看到“合法性、正确性、性能”是三道连续但不同的设计关卡。

\discussionpoint{跨设备配置应是一组经验证的契约。}
即使某台设备允许 1024 线程块，也不能据此否定题设设备的 512 上限，更不能把一次运行经验推广到全部部署节点。生产代码可以查询 \code{cudaDeviceProp} 和函数属性，按计算能力、数据类型与问题形状选择经过验证的配置族，或把选择交给具有调优记录的库；查询结果仍须与内核的动态共享内存和寄存器需求联合检查。持续集成应至少编译主要目标 SM，运行合法边界和故意超限的负测试，并确认错误从异步调用链传到测试结果。这样“曾经在一台 GPU 上运行”才会升级为可审计、可移植的部署证据。
\variantheading 改为 $32\times16$ 线程块后每块正好 512 线程，覆盖 $1024^2$ 输出需 $32\times64=2048$ 个块。
\practiceheading 每次启动后立即检查 \code{cudaGetLastError}，并在测试同步点检查执行错误；部署时查询 \code{cudaDeviceProp.maxThreadsPerBlock} 而非硬编码设备能力。
\sourcecheck{https://docs.nvidia.com/cuda/cuda-c-programming-guide/}{NVIDIA CUDA C++ Programming Guide 的每块最大线程数与内核启动错误规则；配置乘积独立复核，置信度高}
\end{solutionexercise}
